\documentclass[11pt]{article}

% ---- theme variables (passed via pandoc -V) ----
\newcommand{\ThemeName}{System Design Problems — Foundational}
\newcommand{\ThemeKicker}{Design Track}
\newcommand{\ThemeNumber}{I}

\usepackage{interviewstyle}
\setthemecolor{db2777}{f472b6}

% pandoc-required plumbing
\providecommand{\tightlist}{\setlength{\itemsep}{1pt}\setlength{\parskip}{0pt}}
\setlength{\emergencystretch}{3em}
\providecommand{\pandocbounded}[1]{#1}

% syntax-highlighting macros emitted by pandoc (skylighting)
\usepackage{color}
\usepackage{fancyvrb}
\newcommand{\VerbBar}{|}
\newcommand{\VERB}{\Verb[commandchars=\\\{\}]}
\DefineVerbatimEnvironment{Highlighting}{Verbatim}{commandchars=\\\{\}}
% Add ',fontsize=\small' for more characters per line
\usepackage{framed}
\definecolor{shadecolor}{RGB}{35,38,41}
\newenvironment{Shaded}{\begin{snugshade}}{\end{snugshade}}
\newcommand{\AlertTok}[1]{\textcolor[rgb]{0.58,0.85,0.30}{\textbf{\colorbox[rgb]{0.30,0.12,0.14}{#1}}}}
\newcommand{\AnnotationTok}[1]{\textcolor[rgb]{0.25,0.50,0.35}{#1}}
\newcommand{\AttributeTok}[1]{\textcolor[rgb]{0.16,0.50,0.73}{#1}}
\newcommand{\BaseNTok}[1]{\textcolor[rgb]{0.96,0.45,0.00}{#1}}
\newcommand{\BuiltInTok}[1]{\textcolor[rgb]{0.50,0.55,0.55}{#1}}
\newcommand{\CharTok}[1]{\textcolor[rgb]{0.24,0.68,0.91}{#1}}
\newcommand{\CommentTok}[1]{\textcolor[rgb]{0.48,0.49,0.49}{#1}}
\newcommand{\CommentVarTok}[1]{\textcolor[rgb]{0.50,0.55,0.55}{#1}}
\newcommand{\ConstantTok}[1]{\textcolor[rgb]{0.15,0.68,0.68}{\textbf{#1}}}
\newcommand{\ControlFlowTok}[1]{\textcolor[rgb]{0.99,0.74,0.29}{\textbf{#1}}}
\newcommand{\DataTypeTok}[1]{\textcolor[rgb]{0.16,0.50,0.73}{#1}}
\newcommand{\DecValTok}[1]{\textcolor[rgb]{0.96,0.45,0.00}{#1}}
\newcommand{\DocumentationTok}[1]{\textcolor[rgb]{0.64,0.20,0.25}{#1}}
\newcommand{\ErrorTok}[1]{\textcolor[rgb]{0.85,0.27,0.33}{\underline{#1}}}
\newcommand{\ExtensionTok}[1]{\textcolor[rgb]{0.00,0.60,1.00}{\textbf{#1}}}
\newcommand{\FloatTok}[1]{\textcolor[rgb]{0.96,0.45,0.00}{#1}}
\newcommand{\FunctionTok}[1]{\textcolor[rgb]{0.56,0.27,0.68}{#1}}
\newcommand{\ImportTok}[1]{\textcolor[rgb]{0.15,0.68,0.38}{#1}}
\newcommand{\InformationTok}[1]{\textcolor[rgb]{0.77,0.36,0.00}{#1}}
\newcommand{\KeywordTok}[1]{\textcolor[rgb]{0.81,0.81,0.76}{\textbf{#1}}}
\newcommand{\NormalTok}[1]{\textcolor[rgb]{0.81,0.81,0.76}{#1}}
\newcommand{\OperatorTok}[1]{\textcolor[rgb]{0.81,0.81,0.76}{#1}}
\newcommand{\OtherTok}[1]{\textcolor[rgb]{0.15,0.68,0.38}{#1}}
\newcommand{\PreprocessorTok}[1]{\textcolor[rgb]{0.15,0.68,0.38}{#1}}
\newcommand{\RegionMarkerTok}[1]{\textcolor[rgb]{0.16,0.50,0.73}{\colorbox[rgb]{0.08,0.19,0.26}{#1}}}
\newcommand{\SpecialCharTok}[1]{\textcolor[rgb]{0.24,0.68,0.91}{#1}}
\newcommand{\SpecialStringTok}[1]{\textcolor[rgb]{0.85,0.27,0.33}{#1}}
\newcommand{\StringTok}[1]{\textcolor[rgb]{0.96,0.31,0.31}{#1}}
\newcommand{\VariableTok}[1]{\textcolor[rgb]{0.15,0.68,0.68}{#1}}
\newcommand{\VerbatimStringTok}[1]{\textcolor[rgb]{0.85,0.27,0.33}{#1}}
\newcommand{\WarningTok}[1]{\textcolor[rgb]{0.85,0.27,0.33}{#1}}

\begin{document}

\makecoverpage

\thispagestyle{empty}
{\hypersetup{hidelinks}\tableofcontents}
\clearpage

\begin{quote}
Complete end-to-end worked designs for a 45--60 min FAANG system-design round.
Each section mirrors the structure you should drive in the actual interview.
\end{quote}

\section{1. Design a URL Shortener (TinyURL / bit.ly)}\label{1-design-a-url-shortener-tinyurl--bitly}

\subsection{1. Requirements}\label{1-requirements}

\textbf{Clarifying questions to ask the interviewer:}

\begin{itemize}
\tightlist
\item
  What is the expected scale? (DAU, writes per second, reads per second?)
\item
  Should we support custom aliases (user-defined short codes)?
\item
  What is the required URL lifetime? Do links expire?
\item
  Do we need analytics (click counts, geo, referrer)?
\item
  Should we support link editing or deletion?
\item
  Is high availability or strong consistency more important?
\end{itemize}

\textbf{Functional Requirements:}

\begin{enumerate}
\def\labelenumi{\arabic{enumi}.}
\tightlist
\item
  Given a long URL, return a unique short URL (e.g., \texttt{https://tny.io/aB3kX9}).
\item
  Redirect a short URL to the original long URL.
\item
  (Optional) Custom aliases: user can specify the short code.
\item
  (Optional) URL expiry: links expire after N days.
\item
  (Optional) Analytics: track click count, geography, referrer.
\end{enumerate}

\textbf{Non-Functional Requirements:}

\begin{itemize}
\tightlist
\item
  100 ms P99 redirect latency.
\item
  99.99\% availability (URL shortener is a critical dependency for customers).
\item
  Short codes must be unique and not predictable/enumerable.
\item
  The system is \textbf{read-heavy} (100:1 read-to-write ratio is typical).
\end{itemize}

\subsection{2. Capacity Estimation}\label{2-capacity-estimation}

\textbf{Assumptions:}

\begin{itemize}
\tightlist
\item
  500 million URLs shortened per month (writes).
\item
  100:1 read-to-write ratio \(\rightarrow\) 50 billion redirects per month.
\end{itemize}

\textbf{Write QPS:}

\setcodelang{}

\begin{verbatim}
500,000,000 / (30 × 86,400) ≈ 193 writes/sec  ≈ ~200 writes/sec
Peak (3×):                                      ≈ ~600 writes/sec
\end{verbatim}

\textbf{Read QPS:}

\setcodelang{}

\begin{verbatim}
50,000,000,000 / (30 × 86,400) ≈ 19,290 reads/sec  ≈ ~20,000 reads/sec
Peak (3×):                                           ≈ ~60,000 reads/sec
\end{verbatim}

\textbf{Storage (URLs, 5 years):}

\setcodelang{Python}

\begin{Shaded}
\begin{Highlighting}[]
\DecValTok{500}\ErrorTok{M}\NormalTok{ writes}\OperatorTok{/}\NormalTok{month × }\DecValTok{12}\NormalTok{ months × }\DecValTok{5}\NormalTok{ years }\OperatorTok{=} \DecValTok{30}\NormalTok{ billion URLs}
\NormalTok{Each record ≈ }\DecValTok{500} \BuiltInTok{bytes}\NormalTok{ (}\BuiltInTok{long}\NormalTok{ URL }\OperatorTok{\textasciitilde{}}\DecValTok{400}\NormalTok{ B }\OperatorTok{+}\NormalTok{ metadata }\OperatorTok{\textasciitilde{}}\DecValTok{100}\NormalTok{ B)}
\NormalTok{Total: }\DecValTok{30}\ErrorTok{B}\NormalTok{ × }\DecValTok{500}\NormalTok{ B }\OperatorTok{=} \DecValTok{15}\NormalTok{ TB}
\end{Highlighting}
\end{Shaded}

\textbf{Cache (top 20\% URLs serve 80\% traffic):}

\setcodelang{}

\begin{verbatim}
20,000 reads/sec × 86,400 sec/day = 1.728 billion reads/day
20% of 30B URLs = 6 billion "hot" URLs → not all fit in RAM
Pragmatic: cache last 24 hrs of writes = 500M/30 × 500B ≈ 8.3 GB/day  →  ~10 GB RAM
\end{verbatim}

\textbf{Bandwidth:}

\setcodelang{}

\begin{verbatim}
Inbound (writes):  200/sec × 500 B  ≈ 100 KB/s
Outbound (reads):  20,000/sec × 500 B ≈ 10 MB/s
\end{verbatim}

\subsection{3. API Design}\label{3-api-design}

\setcodelang{Python}

\begin{Shaded}
\begin{Highlighting}[]
\CommentTok{\# Shorten a URL}
\NormalTok{POST }\OperatorTok{/}\NormalTok{api}\OperatorTok{/}\NormalTok{v1}\OperatorTok{/}\NormalTok{shorten}
\NormalTok{Request:}
\NormalTok{  \{}
    \StringTok{"long\_url"}\NormalTok{:   }\StringTok{"https://www.example.com/very/long/path?q=something"}\NormalTok{,}
    \StringTok{"custom\_code"}\NormalTok{: }\StringTok{"my{-}alias"}\NormalTok{,          }\OperatorTok{//}\NormalTok{ optional}
    \StringTok{"expires\_at"}\NormalTok{:  }\StringTok{"2027{-}01{-}01T00:00Z"} \OperatorTok{//}\NormalTok{ optional}
\NormalTok{  \}}
\NormalTok{Response }\DecValTok{201}\NormalTok{:}
\NormalTok{  \{}
    \StringTok{"short\_url"}\NormalTok{: }\StringTok{"https://tny.io/aB3kX9"}\NormalTok{,}
    \StringTok{"short\_code"}\NormalTok{: }\StringTok{"aB3kX9"}\NormalTok{,}
    \StringTok{"expires\_at"}\NormalTok{: }\StringTok{"2027{-}01{-}01T00:00Z"}
\NormalTok{  \}}

\CommentTok{\# Redirect (the critical read path)}
\NormalTok{GET }\OperatorTok{/}\NormalTok{\{short\_code\}}
\NormalTok{Response }\DecValTok{302}\NormalTok{:}
\NormalTok{  Location: https:}\OperatorTok{//}\NormalTok{www.example.com}\OperatorTok{/}\NormalTok{very}\OperatorTok{/}\BuiltInTok{long}\OperatorTok{/}\NormalTok{path?q}\OperatorTok{=}\NormalTok{something}

\CommentTok{\# Analytics (optional)}
\NormalTok{GET }\OperatorTok{/}\NormalTok{api}\OperatorTok{/}\NormalTok{v1}\OperatorTok{/}\NormalTok{stats}\OperatorTok{/}\NormalTok{\{short\_code\}}
\NormalTok{Response }\DecValTok{200}\NormalTok{:}
\NormalTok{  \{}
    \StringTok{"clicks"}\NormalTok{: }\DecValTok{14982}\NormalTok{,}
    \StringTok{"top\_countries"}\NormalTok{: [...],}
    \StringTok{"referrers"}\NormalTok{: [...]}
\NormalTok{  \}}

\CommentTok{\# Delete}
\NormalTok{DELETE }\OperatorTok{/}\NormalTok{api}\OperatorTok{/}\NormalTok{v1}\OperatorTok{/}\NormalTok{shorten}\OperatorTok{/}\NormalTok{\{short\_code\}}
\NormalTok{Response }\DecValTok{204}
\end{Highlighting}
\end{Shaded}

\subsection{4. Data Model}\label{4-data-model}

\textbf{SQL vs NoSQL decision:}

\begin{itemize}
\tightlist
\item
  The primary access pattern is a simple key-value lookup: \texttt{short\_code\ →\ long\_url}.
\item
  No complex joins needed; writes are straightforward inserts.
\item
  \textbf{Choice: NoSQL key-value store (Cassandra or DynamoDB) for the URL table} --- horizontal scale, high availability, and the access pattern is a perfect fit.
\item
  A relational DB (PostgreSQL) is used for user/account data where consistency matters.
\end{itemize}

\textbf{URL Table (Cassandra / DynamoDB):}

\setcodelang{Python}

\begin{Shaded}
\begin{Highlighting}[]
\NormalTok{urls}

\NormalTok{short\_code    TEXT       PRIMARY KEY        }\OperatorTok{{-}{-}} \StringTok{"aB3kX9"}
\NormalTok{long\_url      TEXT       NOT NULL           }\OperatorTok{{-}{-}}\NormalTok{ original URL}
\NormalTok{user\_id       UUID                          }\OperatorTok{{-}{-}}\NormalTok{ creator (nullable }\ControlFlowTok{for}\NormalTok{ anonymous)}
\NormalTok{created\_at    TIMESTAMP}
\NormalTok{expires\_at    TIMESTAMP                     }\OperatorTok{{-}{-}}\NormalTok{ NULL }\OperatorTok{=}\NormalTok{ never expires}
\NormalTok{is\_deleted    BOOLEAN    DEFAULT false}
\end{Highlighting}
\end{Shaded}

\textbf{Users Table (PostgreSQL):}

\setcodelang{}

\begin{verbatim}
users

user_id       UUID       PRIMARY KEY
email         VARCHAR(255) UNIQUE NOT NULL
api_key       VARCHAR(64)  UNIQUE            -- for rate limiting
created_at    TIMESTAMP
\end{verbatim}

\textbf{Analytics Table (write-optimized, e.g., Cassandra or ClickHouse):}

\setcodelang{Python}

\begin{Shaded}
\begin{Highlighting}[]
\NormalTok{click\_events}

\NormalTok{short\_code    TEXT}
\NormalTok{clicked\_at    TIMESTAMP}
\NormalTok{country       TEXT}
\NormalTok{referrer      TEXT}
\NormalTok{user\_agent    TEXT}
\OperatorTok{{-}{-}}\NormalTok{ Partition key: short\_code}\OperatorTok{;}\NormalTok{ Clustering key: clicked\_at DESC}
\end{Highlighting}
\end{Shaded}

\subsection{5. High-Level Architecture}\label{5-high-level-architecture}

\setcodelang{}

\begin{verbatim}
                         ┌─────────────────────────────────────────────────┐
                         │                   Clients                       │
                         │        (Browser / Mobile / API consumers)       │
                         └──────────────────────┬──────────────────────────┘
                                                │ HTTPS
                                     ┌──────────▼──────────┐
                                     │   Load Balancer      │
                                     │  (L7, e.g. Nginx /  │
                                     │   AWS ALB)           │
                                     └──────┬───────┬───────┘
                                            │       │
                               ┌────────────▼──┐ ┌──▼────────────┐
                               │  Write Service │ │ Redirect Svc  │
                               │  (POST /shorten│ │ (GET /{code}) │
                               └───────┬────────┘ └──────┬────────┘
                                       │                  │
                              ┌────────▼──────┐   ┌───────▼────────┐
                              │  Key Gen Svc  │   │  Redis Cache   │
                              │  (KGS — see   │   │  (short_code → │
                              │   Deep Dive)  │   │   long_url)    │
                              └────────┬──────┘   └───────┬────────┘
                                       │                   │ cache miss
                                       │            ┌──────▼─────────┐
                                       └───────────►│  Cassandra DB  │
                                                    │  (URL table)   │
                                                    └──────┬─────────┘
                                                           │
                                                  ┌────────▼────────┐
                                                  │  Kafka (async)  │
                                                  │  Click events   │
                                                  └────────┬────────┘
                                                           │
                                                  ┌────────▼────────┐
                                                  │  Analytics Svc  │
                                                  │  (ClickHouse)   │
                                                  └─────────────────┘
\end{verbatim}

\textbf{Component walkthrough:}

\begin{enumerate}
\def\labelenumi{\arabic{enumi}.}
\tightlist
\item
  \textbf{Client} sends \texttt{POST\ /shorten} or \texttt{GET\ /\{code\}} over HTTPS.
\item
  \textbf{Load Balancer} routes to stateless Write Service or Redirect Service pods.
\item
  \textbf{Redirect Service} checks \textbf{Redis} first (cache hit \(\rightarrow\) 302 immediately). On cache miss, reads Cassandra, populates cache, then redirects.
\item
  \textbf{Write Service} calls \textbf{KGS} (Key Generation Service) for a pre-generated unique code, stores \texttt{(code,\ long\_url)} in Cassandra, and warms the cache.
\item
  \textbf{Kafka} asynchronously receives click events; \textbf{Analytics Service} batch-writes to ClickHouse for reporting.
\end{enumerate}

\subsection{6. Deep Dives}\label{6-deep-dives}

\subsubsection{6.1 Key Generation: How to create short codes?}\label{61-key-generation-how-to-create-short-codes}

\begin{longtable}[]{@{}
  >{\raggedright\arraybackslash}p{(\columnwidth - 6\tabcolsep) * \real{0.1103}}
  >{\raggedright\arraybackslash}p{(\columnwidth - 6\tabcolsep) * \real{0.3278}}
  >{\raggedright\arraybackslash}p{(\columnwidth - 6\tabcolsep) * \real{0.2616}}
  >{\raggedright\arraybackslash}p{(\columnwidth - 6\tabcolsep) * \real{0.3002}}@{}}
\toprule\noalign{}
\rowcolor{acc!16}
\begin{minipage}[b]{\linewidth}\raggedright
Approach
\end{minipage} & \begin{minipage}[b]{\linewidth}\raggedright
How
\end{minipage} & \begin{minipage}[b]{\linewidth}\raggedright
Pros
\end{minipage} & \begin{minipage}[b]{\linewidth}\raggedright
Cons
\end{minipage} \\
\midrule\noalign{}
\endhead
\bottomrule\noalign{}
\endlastfoot
\textbf{MD5 / SHA-256 hash} & Hash long URL, take first 6--8 chars & Deterministic (same URL \(\rightarrow\) same code) & Collisions possible; must check DB on every write \\
\rowcolor{zebra}
\textbf{Counter + Base62} & Global atomic counter \(\rightarrow\) encode in Base62 & No collisions; simple & Single counter = SPOF; reveals ordering/volume \\
\textbf{Pre-generated KGS} & Dedicated service pre-generates \& stores unused keys in a "keys" table & No collision checking at write time; fast & Additional service to maintain; key theft if service crashes \\
\rowcolor{zebra}
\textbf{UUID} & Generate random UUID, take first 8 chars & Simple & High collision probability at scale \\
\end{longtable}

\textbf{Key Decision: KGS (Key Generation Service) is the recommended approach for large scale.}

KGS design:

\begin{itemize}
\tightlist
\item
  KGS pre-generates 6-character Base62 codes (62\^{}6 \(\approx\) 56 billion unique codes).
\item
  Stores them in two tables: \texttt{unused\_keys} and \texttt{used\_keys}.
\item
  Each Write Service instance requests a batch (e.g., 1,000 keys) from KGS on startup, holds them in-memory.
\item
  KGS marks those keys as "in use" atomically. If a Write Service pod crashes, those keys are wasted but never reused --- acceptable loss.
\item
  KGS itself can be replicated with a leader for key assignment.
\end{itemize}

\subsubsection{6.2 Redirect: 301 vs 302?}\label{62-redirect-301-vs-302}

\begin{longtable}[]{@{}
  >{\raggedright\arraybackslash}p{(\columnwidth - 4\tabcolsep) * \real{0.1866}}
  >{\raggedright\arraybackslash}p{(\columnwidth - 4\tabcolsep) * \real{0.3940}}
  >{\raggedright\arraybackslash}p{(\columnwidth - 4\tabcolsep) * \real{0.4194}}@{}}
\toprule\noalign{}
\rowcolor{acc!16}
\begin{minipage}[b]{\linewidth}\raggedright
\end{minipage} & \begin{minipage}[b]{\linewidth}\raggedright
301 Permanent
\end{minipage} & \begin{minipage}[b]{\linewidth}\raggedright
302 Temporary
\end{minipage} \\
\midrule\noalign{}
\endhead
\bottomrule\noalign{}
\endlastfoot
Browser caches? & Yes --- browser won\textquotesingle t call TinyURL again & No --- browser calls TinyURL on every click \\
\rowcolor{zebra}
Analytics accuracy & \textbf{Broken} --- browser bypasses your servers & \textbf{Accurate} --- every click hits your servers \\
Load on TinyURL & Lower after first visit & Higher (every click) \\
\rowcolor{zebra}
\textbf{Recommendation} & Use only if analytics not needed & \textbf{Use 302 for full analytics capability} \\
\end{longtable}

\textbf{Interview takeaway:} Always use 302 unless you explicitly trade analytics for reduced server load.

\subsubsection{6.3 Collision Handling}\label{63-collision-handling}

If using hash-based approach:

\begin{enumerate}
\def\labelenumi{\arabic{enumi}.}
\tightlist
\item
  Compute MD5(long\_url) \(\rightarrow\) take first 7 chars as short\_code.
\item
  Check if \texttt{short\_code} exists in DB.
\item
  If collision: append a salt (e.g., increment integer), re-hash, retry.
\item
  Limit retries to 3; fall back to random generation if exhausted.
\end{enumerate}

With KGS, collisions are impossible --- keys are pre-assigned and marked used atomically.

\subsubsection{6.4 Custom Aliases}\label{64-custom-aliases}

\begin{itemize}
\tightlist
\item
  User provides desired short code (e.g., \texttt{my-brand}).
\item
  Check availability: \texttt{SELECT\ 1\ FROM\ urls\ WHERE\ short\_code\ =\ \textquotesingle{}my-brand\textquotesingle{}}.
\item
  If taken \(\rightarrow\) return 409 Conflict.
\item
  Store with \texttt{is\_custom\ =\ true} flag to distinguish from generated codes.
\item
  Rate-limit custom alias creation per user to prevent squatting.
\end{itemize}

\subsubsection{6.5 URL Expiry}\label{65-url-expiry}

\begin{itemize}
\tightlist
\item
  Store \texttt{expires\_at} timestamp in the URL record.
\item
  On redirect: if \texttt{expires\_at\ \textless{}\ NOW()} \(\rightarrow\) return 404 (or 410 Gone).
\item
  Background cleanup job: delete expired records weekly (keep DB lean).
\item
  Cache entries: set Redis TTL = \texttt{expires\_at\ -\ NOW()} so expired URLs auto-evict.
\end{itemize}

\subsection{7. Bottlenecks \& Scaling}\label{7-bottlenecks--scaling}

\textbf{Read hotspot --- viral URLs:}

\begin{itemize}
\tightlist
\item
  A single URL (e.g., link from a viral tweet) can receive millions of reads/sec.
\item
  Redis handles this well --- single key lookup is O(1) and Redis handles 100k+ ops/sec per node.
\item
  For extreme cases: replicate the specific key across multiple Redis nodes or use local in-process cache (LRU with 1000 entries) in each Redirect Service pod.
\end{itemize}

\textbf{Write scaling:}

\begin{itemize}
\tightlist
\item
  Writes at \textasciitilde200/sec are trivially handled. Cassandra scales horizontally; add nodes as needed.
\item
  KGS: single KGS with a read replica for HA is sufficient at this scale. Key generation is CPU-light.
\end{itemize}

\textbf{DB sharding:}

\begin{itemize}
\tightlist
\item
  Cassandra partitions by \texttt{short\_code} automatically (consistent hashing). No manual sharding needed.
\item
  If using SQL: shard by \texttt{short\_code} modulo N (hash partitioning). Use a consistent hash ring to avoid resharding pain.
\end{itemize}

\textbf{Cache eviction:}

\begin{itemize}
\tightlist
\item
  Redis with LRU policy. Set max memory = 10--20 GB.
\item
  TTL on each entry = min(expires\_at, 24 hrs) to prevent stale data.
\end{itemize}

\textbf{Failure handling:}

\begin{itemize}
\tightlist
\item
  Redis down \(\rightarrow\) fall through to Cassandra (latency increases but system stays up).
\item
  KGS down \(\rightarrow\) Write Service uses its in-memory key batch until KGS recovers; alert if batch exhausted.
\item
  Cassandra node down \(\rightarrow\) Cassandra replication factor 3 (RF=3) ensures reads/writes still succeed with 2 nodes up.
\end{itemize}

\textbf{Geographic distribution:}

\begin{itemize}
\tightlist
\item
  Deploy read replicas of Cassandra in multiple regions.
\item
  Use GeoDNS to route users to nearest Redirect Service cluster.
\item
  URL writes go to a primary region; replicated asynchronously to others.
\end{itemize}

\subsection{8. Trade-offs Summary}\label{8-trade-offs-summary}

\begin{longtable}[]{@{}
  >{\raggedright\arraybackslash}p{(\columnwidth - 6\tabcolsep) * \real{0.1494}}
  >{\raggedright\arraybackslash}p{(\columnwidth - 6\tabcolsep) * \real{0.2988}}
  >{\raggedright\arraybackslash}p{(\columnwidth - 6\tabcolsep) * \real{0.1826}}
  >{\raggedright\arraybackslash}p{(\columnwidth - 6\tabcolsep) * \real{0.3693}}@{}}
\toprule\noalign{}
\rowcolor{acc!16}
\begin{minipage}[b]{\linewidth}\raggedright
Decision
\end{minipage} & \begin{minipage}[b]{\linewidth}\raggedright
Choice Made
\end{minipage} & \begin{minipage}[b]{\linewidth}\raggedright
Alternative
\end{minipage} & \begin{minipage}[b]{\linewidth}\raggedright
Reason
\end{minipage} \\
\midrule\noalign{}
\endhead
\bottomrule\noalign{}
\endlastfoot
Key generation & \textbf{KGS (pre-generated)} & Hash + collision check & No collision risk; predictable latency \\
\rowcolor{zebra}
DB for URLs & \textbf{Cassandra (NoSQL)} & PostgreSQL & Write scalability; simple key-value access pattern \\
Redirect type & \textbf{302 Temporary} & 301 Permanent & Enables accurate analytics \\
\rowcolor{zebra}
Cache & \textbf{Redis (LRU)} & Memcached & Rich data types; TTL support; persistence \\
Analytics pipeline & \textbf{Kafka \(\rightarrow\) ClickHouse (async)} & Sync write to SQL & Decouples write path; handles bursty click events \\
\rowcolor{zebra}
Custom aliases & \textbf{Supported (rate-limited)} & Not supported & Product differentiator; must prevent abuse \\
\end{longtable}

\section{2. Design a Rate Limiter}\label{2-design-a-rate-limiter}

\subsection{1. Requirements}\label{1-requirements-1}

\textbf{Clarifying questions to ask the interviewer:}

\begin{itemize}
\tightlist
\item
  Where should the rate limiter live? Client-side, API gateway, or per-service middleware?
\item
  What are we rate-limiting on? Per user? Per IP? Per API key? Per endpoint?
\item
  What happens when a limit is exceeded --- hard reject (429) or soft throttle (queue)?
\item
  Do we need distributed rate limiting (multiple servers sharing state)?
\item
  What are the latency requirements? Rate limit check must not add significant overhead.
\item
  Do rules need to be configurable at runtime without deployments?
\end{itemize}

\textbf{Functional Requirements:}

\begin{enumerate}
\def\labelenumi{\arabic{enumi}.}
\tightlist
\item
  Limit the number of requests a client (user/IP/API key) can make in a time window.
\item
  Return HTTP 429 Too Many Requests when limit is exceeded, with \texttt{Retry-After} header.
\item
  Support multiple limits (e.g., 100 req/min AND 1,000 req/hr for same client).
\item
  Rules are configurable without service restarts.
\end{enumerate}

\textbf{Non-Functional Requirements:}

\begin{itemize}
\tightlist
\item
  Rate limit check must add \textless{} 5 ms P99 latency overhead.
\item
  Highly available --- rate limiter failure should fail open (allow traffic) not fail closed (block everything).
\item
  Works across multiple servers (distributed, consistent limits).
\item
  Accurate enough --- exact precision not required; small over-limit acceptable.
\end{itemize}

\subsection{2. Capacity Estimation}\label{2-capacity-estimation-1}

\textbf{Assumptions:}

\begin{itemize}
\tightlist
\item
  10 million DAU, each making \textasciitilde100 requests/day \(\rightarrow\) 1 billion requests/day.
\item
  Rate limit state must be checked for every request.
\end{itemize}

\textbf{QPS to rate limiter:}

\setcodelang{}

\begin{verbatim}
1,000,000,000 / 86,400 ≈ 11,574 req/sec  ≈ ~12,000 checks/sec
Peak (3×):                                ≈ ~36,000 checks/sec
\end{verbatim}

\textbf{Memory for rate limit counters:}

\setcodelang{}

\begin{verbatim}
Strategy: sliding window counter per user per endpoint, 1-hr window.
10M users × 10 endpoints × (counter + timestamp) ≈ 10M × 10 × 20 bytes = 2 GB
→ Fits comfortably in a Redis cluster.
\end{verbatim}

\textbf{Bandwidth:}

\begin{itemize}
\tightlist
\item
  Rate limit check: small --- just a counter lookup and increment.
\item
  Each Redis call \(\approx\) microseconds; no large data transfer.
\end{itemize}

\subsection{3. API Design}\label{3-api-design-1}

The rate limiter is typically middleware, not a standalone API. But if exposing as a service:

\setcodelang{Python}

\begin{Shaded}
\begin{Highlighting}[]
\CommentTok{\# Check and record a request (called by API gateway)}
\NormalTok{POST }\OperatorTok{/}\NormalTok{rate}\OperatorTok{{-}}\NormalTok{limit}\OperatorTok{/}\NormalTok{check}
\NormalTok{Request:}
\NormalTok{  \{}
    \StringTok{"client\_id"}\NormalTok{: }\StringTok{"user:1234"}\NormalTok{,}
    \StringTok{"endpoint"}\NormalTok{:  }\StringTok{"/api/v1/posts"}\NormalTok{,}
    \StringTok{"timestamp"}\NormalTok{: }\DecValTok{1718000000}
\NormalTok{  \}}
\NormalTok{Response }\DecValTok{200}\NormalTok{ (allowed):}
\NormalTok{  \{}
    \StringTok{"allowed"}\NormalTok{: true,}
    \StringTok{"remaining"}\NormalTok{: }\DecValTok{42}\NormalTok{,}
    \StringTok{"reset\_at"}\NormalTok{: }\DecValTok{1718003600}
\NormalTok{  \}}
\NormalTok{Response }\DecValTok{429}\NormalTok{ (throttled):}
\NormalTok{  \{}
    \StringTok{"allowed"}\NormalTok{: false,}
    \StringTok{"retry\_after"}\NormalTok{: }\DecValTok{58}\NormalTok{,       }\OperatorTok{//}\NormalTok{ seconds until window resets}
    \StringTok{"limit"}\NormalTok{: }\DecValTok{100}\NormalTok{,}
    \StringTok{"window"}\NormalTok{: }\StringTok{"1m"}
\NormalTok{  \}}

\CommentTok{\# Update rules (admin)}
\NormalTok{PUT }\OperatorTok{/}\NormalTok{rate}\OperatorTok{{-}}\NormalTok{limit}\OperatorTok{/}\NormalTok{rules}\OperatorTok{/}\NormalTok{\{client\_id\}}
\NormalTok{Request:}
\NormalTok{  \{}
    \StringTok{"rules"}\NormalTok{: [}
\NormalTok{      \{ }\StringTok{"window"}\NormalTok{: }\StringTok{"1m"}\NormalTok{, }\StringTok{"limit"}\NormalTok{: }\DecValTok{100}\NormalTok{ \},}
\NormalTok{      \{ }\StringTok{"window"}\NormalTok{: }\StringTok{"1h"}\NormalTok{, }\StringTok{"limit"}\NormalTok{: }\DecValTok{1000}\NormalTok{ \}}
\NormalTok{    ]}
\NormalTok{  \}}
\end{Highlighting}
\end{Shaded}

\textbf{Response headers to include on every API response:}

\setcodelang{}

\begin{verbatim}
X-RateLimit-Limit:     100
X-RateLimit-Remaining: 42
X-RateLimit-Reset:     1718003600
Retry-After:           58          // only on 429
\end{verbatim}

\subsection{4. Data Model}\label{4-data-model-1}

\textbf{Redis data structures (no SQL needed --- state is ephemeral):}

\textbf{Token Bucket / Fixed Window counter:}

\setcodelang{Python}

\begin{Shaded}
\begin{Highlighting}[]
\NormalTok{Key:   }\StringTok{"rl:}\SpecialCharTok{\{client\_id\}}\StringTok{:}\SpecialCharTok{\{endpoint\}}\StringTok{:}\SpecialCharTok{\{window\_start\}}\StringTok{"}
\NormalTok{Value: integer (request count)}
\NormalTok{TTL:   window duration (e.g., }\DecValTok{60}\ErrorTok{s} \ControlFlowTok{for} \DecValTok{1}\OperatorTok{{-}}\BuiltInTok{min}\NormalTok{ window)}
\end{Highlighting}
\end{Shaded}

\textbf{Sliding Window Log:}

\setcodelang{Python}

\begin{Shaded}
\begin{Highlighting}[]
\NormalTok{Key:   }\StringTok{"rl:}\SpecialCharTok{\{client\_id\}}\StringTok{:}\SpecialCharTok{\{endpoint\}}\StringTok{"}
\NormalTok{Type:  Redis Sorted Set}
\NormalTok{Member: request\_id (UUID)}
\NormalTok{Score:  timestamp (epoch ms)}
\NormalTok{TTL:   window duration}
\end{Highlighting}
\end{Shaded}

\textbf{Rules config (read rarely, cache in-process):}

\setcodelang{Python}

\begin{Shaded}
\begin{Highlighting}[]
\CommentTok{\# Redis Hash or a config store (Consul / etcd)}
\NormalTok{Key:   }\StringTok{"rule:}\SpecialCharTok{\{client\_id\}}\StringTok{"}
\NormalTok{Value: JSON \{ }\StringTok{"rules"}\NormalTok{: [\{}\StringTok{"window"}\NormalTok{:}\StringTok{"1m"}\NormalTok{,}\StringTok{"limit"}\NormalTok{:}\DecValTok{100}\NormalTok{\}, ...] \}}
\end{Highlighting}
\end{Shaded}

\subsection{5. High-Level Architecture}\label{5-high-level-architecture-1}

\setcodelang{}

\begin{verbatim}
           ┌──────────────────────────────────────────────┐
           │                  Client                       │
           └────────────────────┬─────────────────────────┘
                                │ HTTP request
                     ┌──────────▼──────────┐
                     │    API Gateway /     │
                     │    Reverse Proxy     │
                     │  (Nginx / Kong /     │
                     │   AWS API GW)        │
                     └──────┬──────┬────────┘
                            │      │
               ┌────────────▼──┐   │ (if allowed)
               │  Rate Limiter │   │
               │  Middleware   │   │
               │  (sidecar or  │   │
               │   gateway     │   │
               │   plugin)     │   │
               └────────┬──────┘   │
                        │          │
              ┌─────────▼──────┐   │
              │  Redis Cluster │   │
              │  (counters,    │   │
              │   sliding log) │   │
              └────────────────┘   │
                                   ▼
                         ┌─────────────────┐
                         │  Backend API    │
                         │  Services       │
                         └─────────────────┘

Rules flow:
┌──────────────┐       ┌───────────────┐
│ Config Store │──────►│ Rate Limiter  │
│ (etcd/Consul)│       │ (in-process   │
│  (rules)     │       │  cache 30s)   │
└──────────────┘       └───────────────┘
\end{verbatim}

\textbf{Component walkthrough:}

\begin{enumerate}
\def\labelenumi{\arabic{enumi}.}
\tightlist
\item
  Every request hits the \textbf{API Gateway}.
\item
  Gateway invokes \textbf{Rate Limiter Middleware} (same process or sidecar).
\item
  Middleware reads the rule for \texttt{(client\_id,\ endpoint)} from its \textbf{in-process cache} (refreshed every 30s from etcd/Consul).
\item
  Middleware atomically increments/checks the counter in \textbf{Redis}.
\item
  If allowed: forward request to \textbf{Backend Service}.
\item
  If denied: return 429 immediately with \texttt{Retry-After}.
\end{enumerate}

\subsection{6. Deep Dives}\label{6-deep-dives-1}

\subsubsection{6.1 Algorithm Comparison}\label{61-algorithm-comparison}

\begin{longtable}[]{@{}
  >{\raggedright\arraybackslash}p{(\columnwidth - 8\tabcolsep) * \real{0.1188}}
  >{\raggedright\arraybackslash}p{(\columnwidth - 8\tabcolsep) * \real{0.2817}}
  >{\raggedright\arraybackslash}p{(\columnwidth - 8\tabcolsep) * \real{0.1674}}
  >{\raggedright\arraybackslash}p{(\columnwidth - 8\tabcolsep) * \real{0.2647}}
  >{\raggedright\arraybackslash}p{(\columnwidth - 8\tabcolsep) * \real{0.1674}}@{}}
\toprule\noalign{}
\rowcolor{acc!16}
\begin{minipage}[b]{\linewidth}\raggedright
Algorithm
\end{minipage} & \begin{minipage}[b]{\linewidth}\raggedright
How it works
\end{minipage} & \begin{minipage}[b]{\linewidth}\raggedright
Pros
\end{minipage} & \begin{minipage}[b]{\linewidth}\raggedright
Cons
\end{minipage} & \begin{minipage}[b]{\linewidth}\raggedright
Best for
\end{minipage} \\
\midrule\noalign{}
\endhead
\bottomrule\noalign{}
\endlastfoot
\textbf{Fixed Window Counter} & Reset counter every N seconds & Simple; O(1) memory & Boundary burst: 2\(\times\) limit possible at window edge & Simple internal limits \\
\rowcolor{zebra}
\textbf{Sliding Window Log} & Store timestamp of every request in sorted set; count within window & Accurate & O(n) memory per user (n = requests in window) & Low-volume, high-accuracy needs \\
\textbf{Sliding Window Counter} & Blend current + previous window counts weighted by overlap & Near-accurate, O(1) memory & Slight inaccuracy (\textasciitilde0.003\% over-limit) & \textbf{Best general choice} \\
\rowcolor{zebra}
\textbf{Token Bucket} & Bucket of tokens refilled at rate R; request consumes one & Allows bursts up to bucket size & Harder to reason about exact limits; burst may confuse users & APIs that allow bursting \\
\textbf{Leaky Bucket} & Queue requests; process at constant rate & Smooth output rate & Queue introduces latency; bursty requests get delayed & Network traffic shaping \\
\end{longtable}

\textbf{Key Decision: Sliding Window Counter for most API rate limiting. Token Bucket when you want to allow short bursts.}

\subsubsection{6.2 Sliding Window Counter --- Implementation Detail}\label{62-sliding-window-counter--implementation-detail}

\setcodelang{}

\begin{verbatim}
Window = 1 minute.
Current minute bucket: starts at :00
Previous minute bucket: started at :-60

count = prev_bucket_count × (overlap / window) + curr_bucket_count

Example at :45 into the current minute:
overlap = 60 - 45 = 15 seconds (15/60 = 25% of previous window still counts)
count = prev_count × 0.25 + curr_count

If count > limit → reject.
\end{verbatim}

Redis implementation:

\setcodelang{Lua}

\begin{Shaded}
\begin{Highlighting}[]
\CommentTok{{-}{-} Lua script (atomic)}
\KeywordTok{local} \VariableTok{key\_curr} \OperatorTok{=} \ConstantTok{KEYS}\OperatorTok{[}\DecValTok{1}\OperatorTok{]}  \CommentTok{{-}{-} "rl:user1:1718000060"}
\KeywordTok{local} \VariableTok{key\_prev} \OperatorTok{=} \ConstantTok{KEYS}\OperatorTok{[}\DecValTok{2}\OperatorTok{]}  \CommentTok{{-}{-} "rl:user1:1718000000"}
\KeywordTok{local} \VariableTok{limit}    \OperatorTok{=} \FunctionTok{tonumber}\OperatorTok{(}\ConstantTok{ARGV}\OperatorTok{[}\DecValTok{1}\OperatorTok{])}
\KeywordTok{local} \VariableTok{now}      \OperatorTok{=} \FunctionTok{tonumber}\OperatorTok{(}\ConstantTok{ARGV}\OperatorTok{[}\DecValTok{2}\OperatorTok{])}
\KeywordTok{local} \VariableTok{window}   \OperatorTok{=} \FunctionTok{tonumber}\OperatorTok{(}\ConstantTok{ARGV}\OperatorTok{[}\DecValTok{3}\OperatorTok{])}
\KeywordTok{local} \VariableTok{window\_start} \OperatorTok{=} \VariableTok{now} \OperatorTok{{-}} \OperatorTok{(}\VariableTok{now} \OperatorTok{\%} \VariableTok{window}\OperatorTok{)}
\KeywordTok{local} \VariableTok{overlap}  \OperatorTok{=} \VariableTok{window} \OperatorTok{{-}} \OperatorTok{(}\VariableTok{now} \OperatorTok{{-}} \VariableTok{window\_start}\OperatorTok{)}

\KeywordTok{local} \VariableTok{curr\_count} \OperatorTok{=} \FunctionTok{tonumber}\OperatorTok{(}\VariableTok{redis}\OperatorTok{.}\NormalTok{call}\OperatorTok{(}\StringTok{\textquotesingle{}GET\textquotesingle{}}\OperatorTok{,} \VariableTok{key\_curr}\OperatorTok{)} \KeywordTok{or} \DecValTok{0}\OperatorTok{)}
\KeywordTok{local} \VariableTok{prev\_count} \OperatorTok{=} \FunctionTok{tonumber}\OperatorTok{(}\VariableTok{redis}\OperatorTok{.}\NormalTok{call}\OperatorTok{(}\StringTok{\textquotesingle{}GET\textquotesingle{}}\OperatorTok{,} \VariableTok{key\_prev}\OperatorTok{)} \KeywordTok{or} \DecValTok{0}\OperatorTok{)}
\KeywordTok{local} \VariableTok{weighted}   \OperatorTok{=} \VariableTok{prev\_count} \OperatorTok{*} \OperatorTok{(}\VariableTok{overlap} \OperatorTok{/} \VariableTok{window}\OperatorTok{)} \OperatorTok{+} \VariableTok{curr\_count}

\ControlFlowTok{if} \VariableTok{weighted} \OperatorTok{\textgreater{}=} \VariableTok{limit} \ControlFlowTok{then}
  \ControlFlowTok{return} \DecValTok{0}  \CommentTok{{-}{-} rejected}
\ControlFlowTok{end}
\VariableTok{redis}\OperatorTok{.}\NormalTok{call}\OperatorTok{(}\StringTok{\textquotesingle{}INCR\textquotesingle{}}\OperatorTok{,} \VariableTok{key\_curr}\OperatorTok{)}
\VariableTok{redis}\OperatorTok{.}\NormalTok{call}\OperatorTok{(}\StringTok{\textquotesingle{}EXPIRE\textquotesingle{}}\OperatorTok{,} \VariableTok{key\_curr}\OperatorTok{,} \VariableTok{window} \OperatorTok{*} \DecValTok{2}\OperatorTok{)}
\ControlFlowTok{return} \DecValTok{1}  \CommentTok{{-}{-} allowed}
\end{Highlighting}
\end{Shaded}

\subsubsection{6.3 Where to Place the Rate Limiter?}\label{63-where-to-place-the-rate-limiter}

\begin{longtable}[]{@{}
  >{\raggedright\arraybackslash}p{(\columnwidth - 4\tabcolsep) * \real{0.2279}}
  >{\raggedright\arraybackslash}p{(\columnwidth - 4\tabcolsep) * \real{0.3512}}
  >{\raggedright\arraybackslash}p{(\columnwidth - 4\tabcolsep) * \real{0.4208}}@{}}
\toprule\noalign{}
\rowcolor{acc!16}
\begin{minipage}[b]{\linewidth}\raggedright
Placement
\end{minipage} & \begin{minipage}[b]{\linewidth}\raggedright
Pros
\end{minipage} & \begin{minipage}[b]{\linewidth}\raggedright
Cons
\end{minipage} \\
\midrule\noalign{}
\endhead
\bottomrule\noalign{}
\endlastfoot
\textbf{Client-side} & Zero server load & Easily bypassed; untrustworthy \\
\rowcolor{zebra}
\textbf{API Gateway} & Single enforcement point; language-agnostic & Gateway becomes a bottleneck; less context \\
\textbf{Service middleware} & Per-service granularity; business context & Duplicated logic across services; harder to change rules centrally \\
\rowcolor{zebra}
\textbf{Dedicated Rate Limit Service} & Centralized rules; all services call it & Extra network hop; additional failure point \\
\textbf{Recommended: Gateway + Redis} & Rules centralized in Redis; enforced at gateway & Need to handle Redis failure gracefully \\
\end{longtable}

\subsubsection{6.4 Distributed Rate Limiting --- Race Conditions}\label{64-distributed-rate-limiting--race-conditions}

\textbf{Problem:} Two requests arrive simultaneously on different servers. Both read count=99 (limit=100), both increment to 100, both allow. Result: 2 requests slip through over limit.

\textbf{Solutions:}

\begin{enumerate}
\def\labelenumi{\arabic{enumi}.}
\item
  \textbf{Redis Lua scripts (recommended):} Atomic read-increment-check in a single Lua script. Redis is single-threaded; no race condition possible.
\item
  \textbf{Redis INCR + EXPIRE pattern:}
\end{enumerate}

\setcodelang{Python}

\begin{Shaded}
\begin{Highlighting}[]
\NormalTok{count }\OperatorTok{=}\NormalTok{ INCR(key)}
\ControlFlowTok{if}\NormalTok{ count }\OperatorTok{==} \DecValTok{1}\NormalTok{: EXPIRE(key, window)  }\OperatorTok{{-}{-}} \BuiltInTok{set}\NormalTok{ TTL only on first request}
\ControlFlowTok{if}\NormalTok{ count }\OperatorTok{\textgreater{}}\NormalTok{ limit: }\ControlFlowTok{return} \DecValTok{429}
\end{Highlighting}
\end{Shaded}

Problem: race between INCR and EXPIRE if server crashes. Mitigate with \texttt{SET\ key\ 0\ EX\ window\ NX} on first call.

\begin{enumerate}
\def\labelenumi{\arabic{enumi}.}
\setcounter{enumi}{2}
\tightlist
\item
  \textbf{Sticky sessions / consistent hashing:} Route same client\_id to same rate limiter node. Eliminates distributed counting but adds routing complexity and uneven load.
\end{enumerate}

\textbf{Key Decision: Use Redis Lua scripts for atomicity. No locks, no transactions overhead.}

\subsubsection{6.5 Fail-Open vs Fail-Closed}\label{65-fail-open-vs-fail-closed}

\begin{itemize}
\tightlist
\item
  \textbf{Fail-open (recommended):} If Redis is down, allow all requests. Availability \textgreater{} perfect rate limiting.
\item
  \textbf{Fail-closed:} If Redis is down, reject all requests. Too severe for most use cases --- takes down your API.
\item
  Implementation: wrap Redis call in try-catch; on exception, log and allow the request, alert on-call.
\end{itemize}

\subsection{7. Bottlenecks \& Scaling}\label{7-bottlenecks--scaling-1}

\textbf{Redis as single point of failure:}

\begin{itemize}
\tightlist
\item
  Use Redis Cluster (sharded) --- each client\_id maps to a shard via consistent hashing.
\item
  Redis Sentinel for HA (automatic failover).
\item
  Replication lag between primary and replica: acceptable (small over-limit possible).
\end{itemize}

\textbf{Redis latency:}

\begin{itemize}
\tightlist
\item
  Redis P99 \(\approx\) 0.5--1 ms in-datacenter. Total rate limit overhead \textless{} 2 ms.
\item
  If Redis is co-located on same host (Redis sidecar per gateway node): P99 \textless{} 0.2 ms.
\end{itemize}

\textbf{Rule changes at scale:}

\begin{itemize}
\tightlist
\item
  Store rules in etcd/Consul. Rate limiter middleware subscribes to watch-notifications.
\item
  In-process cache with 30s TTL: small window where old rules apply after update.
\end{itemize}

\textbf{Hotspots (one user generating massive traffic):}

\begin{itemize}
\tightlist
\item
  Rate limiter itself handles this well --- a single Redis key gets hammered, but Redis is designed for this.
\item
  Use local in-process counter as first layer (no Redis call for clearly non-bursting clients), then Redis for enforcement.
\end{itemize}

\subsection{8. Trade-offs Summary}\label{8-trade-offs-summary-1}

\begin{longtable}[]{@{}
  >{\raggedright\arraybackslash}p{(\columnwidth - 6\tabcolsep) * \real{0.1217}}
  >{\raggedright\arraybackslash}p{(\columnwidth - 6\tabcolsep) * \real{0.2154}}
  >{\raggedright\arraybackslash}p{(\columnwidth - 6\tabcolsep) * \real{0.1873}}
  >{\raggedright\arraybackslash}p{(\columnwidth - 6\tabcolsep) * \real{0.4757}}@{}}
\toprule\noalign{}
\rowcolor{acc!16}
\begin{minipage}[b]{\linewidth}\raggedright
Decision
\end{minipage} & \begin{minipage}[b]{\linewidth}\raggedright
Choice Made
\end{minipage} & \begin{minipage}[b]{\linewidth}\raggedright
Alternative
\end{minipage} & \begin{minipage}[b]{\linewidth}\raggedright
Reason
\end{minipage} \\
\midrule\noalign{}
\endhead
\bottomrule\noalign{}
\endlastfoot
Algorithm & \textbf{Sliding Window Counter} & Token Bucket & Near-exact accuracy; O(1) memory; prevents boundary bursts \\
\rowcolor{zebra}
State store & \textbf{Redis Cluster} & In-process memory & Distributed consistency across servers; fast enough \\
Placement & \textbf{API Gateway middleware} & Dedicated service & Avoids extra network hop; centralizes enforcement \\
\rowcolor{zebra}
Atomicity & \textbf{Redis Lua scripts} & Redis WATCH/MULTI & Simpler; no retry logic needed \\
Failure mode & \textbf{Fail-open} & Fail-closed & Availability \textgreater{} perfect accuracy for most systems \\
\rowcolor{zebra}
Rules storage & \textbf{etcd + in-process cache} & DB query per request & Rules rarely change; in-process cache eliminates lookup overhead \\
\end{longtable}

\section{3. Design a News Feed (Facebook/Twitter Timeline)}\label{3-design-a-news-feed-facebooktwitter-timeline}

\subsection{1. Requirements}\label{1-requirements-2}

\textbf{Clarifying questions to ask the interviewer:}

\begin{itemize}
\tightlist
\item
  Is this a "friends" model (Facebook --- bidirectional) or "followers" model (Twitter --- unidirectional)?
\item
  What is the feed ordering --- chronological or ranked by relevance?
\item
  What\textquotesingle s the expected social graph size? Max followers per user?
\item
  Do we support stories (ephemeral) vs. posts (permanent)?
\item
  Is the feed pre-generated (push) or assembled on request (pull)?
\item
  Do we need to support ads injection in the feed?
\end{itemize}

\textbf{Functional Requirements:}

\begin{enumerate}
\def\labelenumi{\arabic{enumi}.}
\tightlist
\item
  User can create a post (text, image, video link).
\item
  User sees a feed of posts from people they follow, ordered by recency or ranking.
\item
  Feed updates in near-real-time (within seconds of post creation).
\item
  Pagination: load N posts, scroll for more.
\item
  Users can like, comment, share (out of scope for this design; focus on feed generation).
\end{enumerate}

\textbf{Non-Functional Requirements:}

\begin{itemize}
\tightlist
\item
  500 million DAU (Facebook scale).
\item
  Feed load latency \textless{} 200 ms.
\item
  Post creation propagates to followers within 5 seconds (eventual consistency is OK).
\item
  Users can have up to 5,000 friends (Facebook) or millions of followers (Twitter celebrities).
\end{itemize}

\subsection{2. Capacity Estimation}\label{2-capacity-estimation-2}

\textbf{Assumptions:}

\begin{itemize}
\tightlist
\item
  500 million DAU.
\item
  Each user views feed 5 times/day \(\rightarrow\) 2.5 billion feed reads/day.
\item
  Each user creates 1 post every 5 days \(\rightarrow\) 100 million new posts/day.
\item
  Average friends/followers: 300.
\end{itemize}

\textbf{Post write QPS:}

\setcodelang{}

\begin{verbatim}
100,000,000 / 86,400 ≈ 1,157 writes/sec  ≈ ~1,200 posts/sec
Peak (5×):                                ≈ ~6,000 posts/sec
\end{verbatim}

\textbf{Feed read QPS:}

\setcodelang{}

\begin{verbatim}
2,500,000,000 / 86,400 ≈ 28,935 reads/sec  ≈ ~29,000 reads/sec
Peak (3×):                                  ≈ ~87,000 reads/sec
\end{verbatim}

\textbf{Fanout writes (push model --- writing to all follower feeds on post):}

\setcodelang{}

\begin{verbatim}
1,200 posts/sec × 300 avg followers = 360,000 feed writes/sec
Celebrity (10M followers): 1 post → 10M writes (burst)
\end{verbatim}

\textbf{Storage:}

\setcodelang{}

\begin{verbatim}
Posts: 100M/day × 1 KB (text) = 100 GB/day text
       Images/videos: separate blob store (S3)
Feed cache: 500M users × 20 cached posts × 1 KB = 10 TB → too big for RAM
  → Cache only active users (10% daily active at peak):
    50M users × 20 posts × 1 KB = 1 TB → feasible with Redis cluster
\end{verbatim}

\subsection{3. API Design}\label{3-api-design-2}

\setcodelang{Python}

\begin{Shaded}
\begin{Highlighting}[]
\CommentTok{\# Create a post}
\NormalTok{POST }\OperatorTok{/}\NormalTok{api}\OperatorTok{/}\NormalTok{v1}\OperatorTok{/}\NormalTok{posts}
\NormalTok{Auth: Bearer token}
\NormalTok{Request:}
\NormalTok{  \{}
    \StringTok{"content"}\NormalTok{:    }\StringTok{"Hello world!"}\NormalTok{,}
    \StringTok{"media\_urls"}\NormalTok{: [}\StringTok{"https://cdn.example.com/img/abc.jpg"}\NormalTok{],}
    \StringTok{"visibility"}\NormalTok{: }\StringTok{"friends"}   \OperatorTok{//} \KeywordTok{or} \StringTok{"public"}\NormalTok{, }\StringTok{"private"}
\NormalTok{  \}}
\NormalTok{Response }\DecValTok{201}\NormalTok{:}
\NormalTok{  \{}
    \StringTok{"post\_id"}\NormalTok{:    }\StringTok{"p\_9x3kA2"}\NormalTok{,}
    \StringTok{"created\_at"}\NormalTok{: }\StringTok{"2026{-}06{-}14T10:00:00Z"}
\NormalTok{  \}}

\CommentTok{\# Get news feed}
\NormalTok{GET }\OperatorTok{/}\NormalTok{api}\OperatorTok{/}\NormalTok{v1}\OperatorTok{/}\NormalTok{feed?limit}\OperatorTok{=}\DecValTok{20}\OperatorTok{\&}\NormalTok{cursor}\OperatorTok{=\textless{}}\NormalTok{opaque\_cursor}\OperatorTok{\textgreater{}}
\NormalTok{Auth: Bearer token}
\NormalTok{Response }\DecValTok{200}\NormalTok{:}
\NormalTok{  \{}
    \StringTok{"posts"}\NormalTok{: [}
\NormalTok{      \{}
        \StringTok{"post\_id"}\NormalTok{:    }\StringTok{"p\_9x3kA2"}\NormalTok{,}
        \StringTok{"author\_id"}\NormalTok{:  }\StringTok{"u\_4f2k"}\NormalTok{,}
        \StringTok{"author\_name"}\NormalTok{:}\StringTok{"Alice"}\NormalTok{,}
        \StringTok{"content"}\NormalTok{:    }\StringTok{"Hello world!"}\NormalTok{,}
        \StringTok{"media\_urls"}\NormalTok{: [...],}
        \StringTok{"like\_count"}\NormalTok{: }\DecValTok{142}\NormalTok{,}
        \StringTok{"comment\_count"}\NormalTok{: }\DecValTok{7}\NormalTok{,}
        \StringTok{"created\_at"}\NormalTok{: }\StringTok{"2026{-}06{-}14T10:00:00Z"}
\NormalTok{      \},}
\NormalTok{      ...}
\NormalTok{    ],}
    \StringTok{"next\_cursor"}\NormalTok{: }\StringTok{"\textless{}opaque\_cursor\textgreater{}"}
\NormalTok{  \}}

\CommentTok{\# Get a user\textquotesingle{}s profile posts}
\NormalTok{GET }\OperatorTok{/}\NormalTok{api}\OperatorTok{/}\NormalTok{v1}\OperatorTok{/}\NormalTok{users}\OperatorTok{/}\NormalTok{\{user\_id\}}\OperatorTok{/}\NormalTok{posts?limit}\OperatorTok{=}\DecValTok{20}\OperatorTok{\&}\NormalTok{cursor}\OperatorTok{=\textless{}}\NormalTok{cursor}\OperatorTok{\textgreater{}}
\NormalTok{Response }\DecValTok{200}\NormalTok{: \{ }\StringTok{"posts"}\NormalTok{: [...], }\StringTok{"next\_cursor"}\NormalTok{: }\StringTok{"..."}\NormalTok{ \}}
\end{Highlighting}
\end{Shaded}

\subsection{4. Data Model}\label{4-data-model-2}

\textbf{SQL vs NoSQL:}

\begin{itemize}
\tightlist
\item
  \textbf{Posts table:} Write-once, read-many. High read QPS. Shard by \texttt{author\_id}. \(\rightarrow\) NoSQL (Cassandra) or sharded PostgreSQL.
\item
  \textbf{Social graph (follows):} Relationship lookups ("who does user X follow?"). \(\rightarrow\) Graph DB (Neo4j) or sharded SQL; at Facebook scale: custom TAO (graph store).
\item
  \textbf{Feed cache:} Redis sorted sets (score = post timestamp).
\item
  \textbf{User data:} Relational (PostgreSQL) --- low write volume, complex queries.
\end{itemize}

\textbf{Posts Table (Cassandra --- partition by author, cluster by time):}

\setcodelang{}

\begin{verbatim}
posts

author_id   UUID        PARTITION KEY
post_id     TIMEUUID    CLUSTERING KEY (DESC)   -- newest first
content     TEXT
media_urls  LIST<TEXT>
like_count  COUNTER
created_at  TIMESTAMP
\end{verbatim}

\textbf{Social Graph Table (Cassandra --- two tables for O(1) lookups):}

\setcodelang{}

\begin{verbatim}
follows_by_follower          -- "who does user X follow?"

follower_id  UUID  PARTITION KEY
followee_id  UUID  CLUSTERING KEY

follows_by_followee          -- "who follows user X?"

followee_id  UUID  PARTITION KEY
follower_id  UUID  CLUSTERING KEY
\end{verbatim}

\textbf{Feed Cache (Redis Sorted Set per user):}

\setcodelang{Python}

\begin{Shaded}
\begin{Highlighting}[]
\NormalTok{Key:     }\StringTok{"feed:}\SpecialCharTok{\{user\_id\}}\StringTok{"}
\NormalTok{Members: post\_id}
\NormalTok{Scores:  timestamp (epoch ms) — higher }\OperatorTok{=}\NormalTok{ newer}
\NormalTok{TTL:     }\DecValTok{48}\NormalTok{ hours (evict inactive users)}
\NormalTok{Max size: }\DecValTok{800}\NormalTok{ posts per user (trim older entries)}
\end{Highlighting}
\end{Shaded}

\subsection{5. High-Level Architecture}\label{5-high-level-architecture-2}

\setcodelang{}

\begin{verbatim}
            ┌──────────────────────────────────────────────────┐
            │                   Client                          │
            │           (iOS / Android / Web)                   │
            └────────────────────┬─────────────────────────────┘
                                 │ HTTPS
                      ┌──────────▼──────────┐
                      │    Load Balancer     │
                      └──────┬──────┬────────┘
                             │      │
               ┌─────────────▼─┐  ┌─▼──────────────┐
               │  Post Service │  │  Feed Service  │
               │  (create post)│  │  (read feed)   │
               └──────┬────────┘  └──────┬─────────┘
                      │                  │
          ┌───────────▼──────┐    ┌──────▼──────────┐
          │   Kafka Topic:   │    │   Redis Cluster  │
          │   "new-posts"    │    │  (feed cache,   │
          └───────┬──────────┘    │   sorted sets)   │
                  │               └──────┬────────────┘
          ┌───────▼──────────┐           │ cache miss
          │   Fanout Service │           ▼
          │  (Notification   │   ┌───────────────────┐
          │   Workers)       │   │  Cassandra (posts)│
          └───────┬──────────┘   └───────────────────┘
                  │
          ┌───────▼──────────┐
          │  Social Graph    │    ┌───────────────────┐
          │  (Cassandra /    │    │   CDN / Blob      │
          │   TAO)           │    │   Store (S3)      │
          └──────────────────┘    └───────────────────┘

Post creation flow:  Post Svc → Cassandra + Kafka
Fanout flow:         Kafka → Fanout Workers → Redis (write to follower feeds)
Feed read flow:      Feed Svc → Redis cache → (miss) Cassandra
\end{verbatim}

\subsection{6. Deep Dives}\label{6-deep-dives-2}

\subsubsection{6.1 Fanout Strategies}\label{61-fanout-strategies}

\begin{longtable}[]{@{}
  >{\raggedright\arraybackslash}p{(\columnwidth - 8\tabcolsep) * \real{0.0979}}
  >{\raggedright\arraybackslash}p{(\columnwidth - 8\tabcolsep) * \real{0.2780}}
  >{\raggedright\arraybackslash}p{(\columnwidth - 8\tabcolsep) * \real{0.1690}}
  >{\raggedright\arraybackslash}p{(\columnwidth - 8\tabcolsep) * \real{0.2860}}
  >{\raggedright\arraybackslash}p{(\columnwidth - 8\tabcolsep) * \real{0.1690}}@{}}
\toprule\noalign{}
\rowcolor{acc!16}
\begin{minipage}[b]{\linewidth}\raggedright
Strategy
\end{minipage} & \begin{minipage}[b]{\linewidth}\raggedright
How
\end{minipage} & \begin{minipage}[b]{\linewidth}\raggedright
Pros
\end{minipage} & \begin{minipage}[b]{\linewidth}\raggedright
Cons
\end{minipage} & \begin{minipage}[b]{\linewidth}\raggedright
Best for
\end{minipage} \\
\midrule\noalign{}
\endhead
\bottomrule\noalign{}
\endlastfoot
\textbf{Fanout-on-Write (Push)} & On post creation, write post\_id to each follower\textquotesingle s feed cache & Feed reads O(1) --- just read from cache & Write amplification: 1 post \(\rightarrow\) N writes; celebrity (10M followers) = 10M cache writes & Average users \\
\rowcolor{zebra}
\textbf{Fanout-on-Read (Pull)} & On feed request, fetch posts from each followee, merge, sort & No write amplification & Feed read O(N followees) --- slow for users following many people & Celebrity accounts, inactive followers \\
\textbf{Hybrid (recommended)} & Push to regular followers; pull for celebrities at read time; don\textquotesingle t push to inactive users & Balances read and write load & More complex logic & \textbf{Large-scale production systems} \\
\end{longtable}

\textbf{Key Decision: Hybrid fanout.}

Implementation:

\begin{itemize}
\tightlist
\item
  Define "celebrity" as users with \textgreater{} 1M followers.
\item
  On post write: fanout-on-write for non-celebrities; for celebrities, mark the post in a "celebrity posts" index only.
\item
  On feed read: read user\textquotesingle s pre-built feed from Redis (push-mode posts) + fetch latest celebrity posts from celebrity index + merge and deduplicate.
\item
  Skip fanout for users who haven\textquotesingle t been active in 7 days (check last\_active timestamp).
\end{itemize}

\subsubsection{6.2 Celebrity Problem (Hotspot)}\label{62-celebrity-problem-hotspot}

A celebrity with 10 million followers posts \(\rightarrow\) need 10 million cache writes within seconds.

\textbf{Mitigation:}

\begin{enumerate}
\def\labelenumi{\arabic{enumi}.}
\tightlist
\item
  \textbf{Async fanout workers:} Put the fanout job on Kafka with 10,000 partitions. Workers consume in parallel --- 10M writes spread across workers.
\item
  \textbf{Pull for celebrity content:} Don\textquotesingle t write celebrity posts to follower feeds. At read time, the top-K celebrities a user follows are fetched separately and injected into the feed.
\item
  \textbf{Pre-fetch:} For celebrities, replicate their latest posts to a separate, highly-cached "celebrity feed" shard. All followers read from the same shard (CDN-like effect).
\end{enumerate}

\subsubsection{6.3 Feed Ranking}\label{63-feed-ranking}

Simple timeline: sort by \texttt{score\ =\ timestamp}.

ML-ranked feed (Facebook/Instagram model):

\setcodelang{Python}

\begin{Shaded}
\begin{Highlighting}[]
\NormalTok{score }\OperatorTok{=}\NormalTok{ f(recency, engagement\_velocity, relationship\_strength, content\_type)}

\NormalTok{engagement\_velocity }\OperatorTok{=}\NormalTok{ likes\_per\_hour }\OperatorTok{/}\NormalTok{ avg\_likes\_per\_hour\_for\_author}
\NormalTok{relationship\_strength }\OperatorTok{=}\NormalTok{ interaction\_history(viewer, author)}
\end{Highlighting}
\end{Shaded}

\textbf{Interview takeaway:} In the interview, acknowledge ranking but defer details. Say: "We\textquotesingle d train an ML model offline, serve predictions via a Feature Store + low-latency inference service. The feed service re-ranks the top 200 candidates before returning the top 20."

\subsubsection{6.4 Pagination}\label{64-pagination}

\textbf{Cursor-based pagination (not offset):}

\begin{itemize}
\tightlist
\item
  Offset pagination (\texttt{LIMIT\ 20\ OFFSET\ 100}) breaks when new posts are inserted --- items shift, causing duplicates/skips.
\item
  Cursor = encoded last-seen post timestamp + post\_id.
\item
  Redis: \texttt{ZRANGEBYSCORE\ feed:\{user\_id\}\ -inf\ \textless{}cursor\_score\textgreater{}\ LIMIT\ 0\ 20}.
\item
  Cursor makes the feed stable even as new posts arrive.
\end{itemize}

\subsection{7. Bottlenecks \& Scaling}\label{7-bottlenecks--scaling-2}

\textbf{Feed cache size:}

\begin{itemize}
\tightlist
\item
  10 TB for active users --- run 50+ Redis nodes (200 GB each) in a cluster.
\item
  Use consistent hashing for user \(\rightarrow\) Redis shard mapping.
\item
  Evict feeds of users inactive for 7 days to save memory.
\end{itemize}

\textbf{Cassandra post storage:}

\begin{itemize}
\tightlist
\item
  Partition by \texttt{author\_id} \(\rightarrow\) hot partition if one author posts very frequently.
\item
  Mitigate: add \texttt{bucket} to partition key: \texttt{(author\_id,\ bucket)} where \texttt{bucket\ =\ created\_at\ /\ 30\ days}.
\end{itemize}

\textbf{Social graph at scale:}

\begin{itemize}
\tightlist
\item
  Facebook built TAO (The Associations and Objects) --- a custom graph store on top of MySQL + Memcache.
\item
  At smaller scale: Cassandra \texttt{follows\_by\_follower} + \texttt{follows\_by\_followee} tables work well.
\end{itemize}

\textbf{Read path latency breakdown (target \textless{} 200 ms):}

\setcodelang{}

\begin{verbatim}
Redis feed read:          5 ms
Celebrity post fetch:    20 ms (parallel)
User profile hydration:  10 ms (from user cache)
Ranking re-score:        15 ms
Network + serialization: 10 ms
Total:                  ~60 ms  (well within 200 ms budget)
\end{verbatim}

\subsection{8. Trade-offs Summary}\label{8-trade-offs-summary-2}

\begin{longtable}[]{@{}
  >{\raggedright\arraybackslash}p{(\columnwidth - 6\tabcolsep) * \real{0.1575}}
  >{\raggedright\arraybackslash}p{(\columnwidth - 6\tabcolsep) * \real{0.2362}}
  >{\raggedright\arraybackslash}p{(\columnwidth - 6\tabcolsep) * \real{0.2012}}
  >{\raggedright\arraybackslash}p{(\columnwidth - 6\tabcolsep) * \real{0.4051}}@{}}
\toprule\noalign{}
\rowcolor{acc!16}
\begin{minipage}[b]{\linewidth}\raggedright
Decision
\end{minipage} & \begin{minipage}[b]{\linewidth}\raggedright
Choice Made
\end{minipage} & \begin{minipage}[b]{\linewidth}\raggedright
Alternative
\end{minipage} & \begin{minipage}[b]{\linewidth}\raggedright
Reason
\end{minipage} \\
\midrule\noalign{}
\endhead
\bottomrule\noalign{}
\endlastfoot
Fanout strategy & \textbf{Hybrid (push+pull)} & Pure push or pure pull & Balances write amplification vs read latency \\
\rowcolor{zebra}
Feed storage & \textbf{Redis sorted sets} & DB query at read time & Sub-10ms reads; DB can\textquotesingle t handle 87K reads/sec directly \\
Feed ordering & \textbf{Reverse-chronological (MVP)} & ML-ranked feed & Simpler; ranked feed is a layer on top \\
\rowcolor{zebra}
Pagination & \textbf{Cursor-based} & Offset-based & Stable under inserts; no duplicate/skip on new posts \\
Post DB & \textbf{Cassandra} & PostgreSQL (sharded) & Write throughput; natural time-series clustering \\
\rowcolor{zebra}
Celebrity handling & \textbf{Pull at read time} & Fanout to all followers & Avoids 10M+ write storms per celebrity post \\
\end{longtable}

\section{4. Design a Chat System (WhatsApp / Messenger)}\label{4-design-a-chat-system-whatsapp--messenger}

\subsection{1. Requirements}\label{1-requirements-3}

\textbf{Clarifying questions to ask the interviewer:}

\begin{itemize}
\tightlist
\item
  1:1 chat only, or also group chat? What is the max group size?
\item
  Do we need real-time delivery (online) and push notifications (offline)?
\item
  Media support (images, video, voice)?
\item
  Do messages need end-to-end encryption?
\item
  Message history: how far back? Searchable?
\item
  Read receipts (sent / delivered / read)?
\item
  Online presence ("last seen")?
\end{itemize}

\textbf{Functional Requirements:}

\begin{enumerate}
\def\labelenumi{\arabic{enumi}.}
\tightlist
\item
  1:1 messaging with real-time delivery.
\item
  Group messaging (up to 500 members per group).
\item
  Online/offline status with "last seen" timestamp.
\item
  Read receipts: sent \(\rightarrow\) delivered \(\rightarrow\) read.
\item
  Push notifications for offline users.
\item
  Message history persisted; searchable.
\item
  Media: images, files (\textless{} 100 MB each).
\end{enumerate}

\textbf{Non-Functional Requirements:}

\begin{itemize}
\tightlist
\item
  500 million DAU; 100 billion messages/day.
\item
  Message delivery latency \textless{} 500 ms (online users).
\item
  High availability: 99.99\%.
\item
  Messages delivered exactly once and in order.
\item
  End-to-end encryption support (Signal Protocol) --- note architecture impact.
\end{itemize}

\subsection{2. Capacity Estimation}\label{2-capacity-estimation-3}

\textbf{Assumptions:}

\begin{itemize}
\tightlist
\item
  500 million DAU.
\item
  Each user sends 40 messages/day \(\rightarrow\) 20 billion messages/day (1:1 + group).
\item
  Average message size: 100 bytes (text).
\item
  Average group has 10 members \(\rightarrow\) group message creates 9 additional deliveries.
\item
  Media: 5\% of messages are media \(\rightarrow\) 1 billion media messages/day.
\end{itemize}

\textbf{Message write QPS:}

\setcodelang{}

\begin{verbatim}
20,000,000,000 / 86,400 ≈ 231,481 msg/sec  ≈ ~230,000 messages/sec
Peak (3×):                                  ≈ ~700,000 messages/sec
\end{verbatim}

\textbf{Storage:}

\setcodelang{}

\begin{verbatim}
Text messages:  20B/day × 100 B         = 2 TB/day
Media:          1B/day × 1 MB avg       = 1 PB/day  → S3/GCS
Metadata:       20B/day × 50 B          = 1 TB/day
Total text+meta ≈ 3 TB/day → 1.1 PB/year (text only)
\end{verbatim}

\textbf{Concurrent WebSocket connections:}

\setcodelang{}

\begin{verbatim}
500M DAU × 30% online simultaneously = 150M concurrent connections
Each connection: 65KB OS overhead
Total: 150M × 65KB ≈ 9.75 TB RAM just for connections
→ Need ~100,000 chat servers @ 1,500 connections each
(In practice: use event-driven servers; Nginx/Go handles 100K connections per server)
→ 1,500 servers (100K connections each)
\end{verbatim}

\subsection{3. API Design}\label{3-api-design-3}

\setcodelang{}

\begin{verbatim}
# WebSocket connection (real-time channel)
ws://chat.example.com/ws?token=<auth_token>

# WebSocket message format (client → server)
{
  "type":      "message",
  "msg_id":    "client-generated-uuid",   // for deduplication
  "to":        "user:5678",               // or "group:9999"
  "content":   "Hey!",
  "media_url": null,
  "timestamp": 1718000000000             // client time ms
}

# WebSocket message format (server → client, push delivery)
{
  "type":        "message",
  "msg_id":      "server-assigned-uuid",
  "from":        "user:1234",
  "to":          "user:5678",
  "content":     "Hey!",
  "server_ts":   1718000000050,          // server timestamp (canonical ordering)
  "seq":         100034                  // sequence number in conversation
}

# ACK (client → server, delivery acknowledgement)
{ "type": "ack", "msg_id": "...", "ack_type": "delivered" }
{ "type": "ack", "msg_id": "...", "ack_type": "read" }

# REST: Send message (fallback for HTTP clients)
POST /api/v1/messages
Request:  { "to": "user:5678", "content": "Hey!" }
Response: { "msg_id": "...", "server_ts": 1718000000050 }

# REST: Fetch message history
GET /api/v1/conversations/{conversation_id}/messages?before=<msg_id>&limit=50
Response: { "messages": [...], "has_more": true }

# REST: Upload media
POST /api/v1/media/upload
Response: { "media_url": "https://cdn.example.com/m/abc.jpg", "expires_in": 86400 }

# REST: Get online status
GET /api/v1/users/{user_id}/presence
Response: { "online": false, "last_seen": "2026-06-14T09:55:00Z" }
\end{verbatim}

\subsection{4. Data Model}\label{4-data-model-3}

\textbf{Storage choices:}

\begin{itemize}
\tightlist
\item
  \textbf{Messages:} HBase / Cassandra --- write-heavy time-series; partition by conversation, cluster by time.
\item
  \textbf{Conversations/Groups:} Cassandra or DynamoDB.
\item
  \textbf{User data:} PostgreSQL (low write, complex queries).
\item
  \textbf{Presence (online status):} Redis (ephemeral, TTL-based).
\item
  \textbf{Message queue for offline delivery:} Kafka or internal queues.
\item
  \textbf{Media:} S3 / GCS (blob).
\end{itemize}

\textbf{Messages Table (HBase / Cassandra):}

\setcodelang{Python}

\begin{Shaded}
\begin{Highlighting}[]
\NormalTok{messages}

\NormalTok{conversation\_id   TEXT       PARTITION KEY    }\OperatorTok{{-}{-}} \StringTok{"conv:}\SpecialCharTok{\{user1\}}\StringTok{\_}\SpecialCharTok{\{user2\}}\StringTok{"} \KeywordTok{or} \StringTok{"group:}\SpecialCharTok{\{id\}}\StringTok{"}
\NormalTok{seq\_num           BIGINT     CLUSTERING KEY DESC  }\OperatorTok{{-}{-}}\NormalTok{ monotonically increasing per conversation}
\NormalTok{msg\_id            UUID       }\OperatorTok{{-}{-}}\NormalTok{ globally unique}
\NormalTok{sender\_id         UUID}
\NormalTok{content           TEXT}
\NormalTok{media\_url         TEXT}
\NormalTok{msg\_type          TEXT       }\OperatorTok{{-}{-}} \StringTok{"text"}\NormalTok{, }\StringTok{"image"}\NormalTok{, }\StringTok{"video"}\NormalTok{, }\StringTok{"system"}
\NormalTok{status            TEXT       }\OperatorTok{{-}{-}} \StringTok{"sent"}\NormalTok{, }\StringTok{"delivered"}\NormalTok{, }\StringTok{"read"}
\NormalTok{created\_at        TIMESTAMP}
\end{Highlighting}
\end{Shaded}

\textbf{Conversations Table:}

\setcodelang{}

\begin{verbatim}
conversations

conversation_id   TEXT    PRIMARY KEY
participants      LIST<UUID>
group_name        TEXT    (null for 1:1)
last_msg_id       UUID
last_msg_at       TIMESTAMP
\end{verbatim}

\textbf{User Conversations Index (reverse lookup):}

\setcodelang{}

\begin{verbatim}
user_conversations   -- "what conversations is user X in?"

user_id              UUID    PARTITION KEY
conversation_id      TEXT    CLUSTERING KEY
last_msg_at          TIMESTAMP  (for sorting)
unread_count         INT
\end{verbatim}

\textbf{Presence Table (Redis):}

\setcodelang{Python}

\begin{Shaded}
\begin{Highlighting}[]
\NormalTok{Key:   }\StringTok{"presence:}\SpecialCharTok{\{user\_id\}}\StringTok{"}
\NormalTok{Value: \{ }\StringTok{"online"}\NormalTok{: true, }\StringTok{"last\_heartbeat"}\NormalTok{: }\DecValTok{1718000000000}\NormalTok{ \}}
\NormalTok{TTL:   }\DecValTok{30}\NormalTok{ seconds (refreshed by heartbeat every }\DecValTok{10}\ErrorTok{s}\NormalTok{)}
\end{Highlighting}
\end{Shaded}

\subsection{5. High-Level Architecture}\label{5-high-level-architecture-3}

\setcodelang{}

\begin{verbatim}
                    ┌──────────────────────────────────────────────┐
                    │               Clients                         │
                    │     (iOS / Android / Web / Desktop)           │
                    └─────┬──────────────┬───────────────┬──────────┘
                          │ WebSocket    │ REST          │ REST
                          │             │(history/media)│(auth/profile)
               ┌──────────▼──────────┐  │               │
               │  WebSocket LB        │  │               │
               │  (L4, sticky by     │  │               │
               │   connection)        │  │               │
               └──────────┬──────────┘  │               │
                          │             │               │
         ┌────────────────▼──┐          │               │
         │  Chat Servers      │          │               │
         │  (WebSocket mgmt) │          │               │
         │  [Stateful —      │          │               │
         │   hold connections]│         │               │
         └─────┬──────┬──────┘          │               │
               │      │                 │               │
  ┌────────────▼──┐ ┌──▼────────────┐   │               │
  │    Kafka      │ │  Redis        │   │               │
  │  (message    │ │  (presence,   │   │               │
  │   routing,   │ │   connection  │   │               │
  │   delivery   │ │   registry)   │   │               │
  │   queue)     │ └───────────────┘   │               │
  └────────┬─────┘                     │               │
           │                           │               │
  ┌────────▼──────────┐     ┌──────────▼──────┐  ┌─────▼──────────┐
  │  Message Service  │     │  History API    │  │  User Service  │
  │  (store, route,   │     │  (REST reads)   │  │  (auth, profile│
  │   push notifs)    │     └────────┬────────┘  └────────────────┘
  └────────┬──────────┘              │
           │                         │
  ┌────────▼───────────────────────▼─────────────┐
  │                 Cassandra / HBase             │
  │              (messages, conversations)        │
  └───────────────────────────────────────────────┘
           │
  ┌────────▼──────────┐         ┌──────────────────┐
  │  Push Notif Svc   │         │   S3 / GCS       │
  │  (APNs, FCM)      │         │   (media blobs)  │
  └───────────────────┘         └──────────────────┘
\end{verbatim}

\textbf{Message delivery flow (1:1, both online):}

\begin{enumerate}
\def\labelenumi{\arabic{enumi}.}
\tightlist
\item
  Sender client sends message over WebSocket to \textbf{Chat Server A}.
\item
  Chat Server A writes message to \textbf{Kafka} (topic: \texttt{messages}).
\item
  \textbf{Message Service} consumes from Kafka, persists to Cassandra.
\item
  Message Service looks up which \textbf{Chat Server} the recipient is connected to (via Redis connection registry: \texttt{"ws\_server:\{user\_id\}"\ →\ "chat-server-B"}).
\item
  Message Service pushes message to \textbf{Chat Server B} (via internal pub/sub).
\item
  Chat Server B delivers to recipient\textquotesingle s WebSocket.
\item
  Recipient client sends \texttt{ack:\ delivered} \(\rightarrow\) Chat Server B \(\rightarrow\) Message Service \(\rightarrow\) Cassandra (update status) \(\rightarrow\) Kafka \(\rightarrow\) Chat Server A \(\rightarrow\) sender gets double-tick.
\end{enumerate}

\subsection{6. Deep Dives}\label{6-deep-dives-3}

\subsubsection{6.1 WebSocket Connection Management}\label{61-websocket-connection-management}

\textbf{Why WebSocket instead of HTTP polling?}

\begin{longtable}[]{@{}
  >{\raggedright\arraybackslash}p{(\columnwidth - 6\tabcolsep) * \real{0.1719}}
  >{\raggedright\arraybackslash}p{(\columnwidth - 6\tabcolsep) * \real{0.3025}}
  >{\raggedright\arraybackslash}p{(\columnwidth - 6\tabcolsep) * \real{0.2005}}
  >{\raggedright\arraybackslash}p{(\columnwidth - 6\tabcolsep) * \real{0.3251}}@{}}
\toprule\noalign{}
\rowcolor{acc!16}
\begin{minipage}[b]{\linewidth}\raggedright
Method
\end{minipage} & \begin{minipage}[b]{\linewidth}\raggedright
How
\end{minipage} & \begin{minipage}[b]{\linewidth}\raggedright
Pros
\end{minipage} & \begin{minipage}[b]{\linewidth}\raggedright
Cons
\end{minipage} \\
\midrule\noalign{}
\endhead
\bottomrule\noalign{}
\endlastfoot
Short polling & Client polls every N seconds & Simple & Wasted requests; high latency (up to N sec) \\
\rowcolor{zebra}
Long polling & Server holds request until message or timeout & Fewer round trips & Still HTTP overhead; server holds threads \\
\textbf{WebSocket} & Full-duplex TCP connection & True real-time; low overhead & Stateful --- servers must track who is connected where \\
\rowcolor{zebra}
SSE (Server-Sent Events) & Server pushes, client uses XHR for sends & Simpler than WS & Unidirectional; need separate channel for sends \\
\end{longtable}

\textbf{Key Decision: WebSocket for all real-time delivery.}

\textbf{Connection registry in Redis:}

\setcodelang{Python}

\begin{Shaded}
\begin{Highlighting}[]
\NormalTok{SET }\StringTok{"conn:}\SpecialCharTok{\{user\_id\}}\StringTok{"} \StringTok{"chat{-}server{-}47"}\NormalTok{ EX }\DecValTok{300}
\end{Highlighting}
\end{Shaded}

\begin{itemize}
\tightlist
\item
  Refreshed every 60s by Chat Server via heartbeat.
\item
  If key expires (server crash), user treated as offline \(\rightarrow\) push notification path.
\end{itemize}

\textbf{Load balancing WebSocket connections:}

\begin{itemize}
\tightlist
\item
  L4 load balancer (not L7) --- connections are long-lived TCP; no HTTP termination per request.
\item
  Use consistent hashing on \texttt{user\_id} for connection routing (optional but reduces cross-server lookups).
\end{itemize}

\subsubsection{6.2 Message Delivery: Online vs Offline}\label{62-message-delivery-online-vs-offline}

\setcodelang{}

\begin{verbatim}
Sender posts message
        │
        ▼
   Is recipient online? (check Redis presence key)
        │
   YES  │  NO
        │         ├──► Persist to Cassandra
        │         └──► Enqueue push notification (APNs/FCM)
        ▼                    │
  Route to recipient's       │
  Chat Server via Redis      ▼
  connection registry   Push Notif Svc sends silent push
        │               → device wakes up, fetches history
        ▼               via REST, renders notification
  Deliver via WebSocket
\end{verbatim}

\textbf{Message deduplication:}

\begin{itemize}
\tightlist
\item
  Client generates a UUID (\texttt{client\_msg\_id}) before sending.
\item
  Server stores \texttt{client\_msg\_id} \(\rightarrow\) \texttt{server\_msg\_id} mapping for 24 hours.
\item
  If retry received with same \texttt{client\_msg\_id}, return existing \texttt{server\_msg\_id} (idempotent).
\end{itemize}

\subsubsection{6.3 Group Chat Delivery}\label{63-group-chat-delivery}

\textbf{Small groups (\(\leq\) 500 members):}

\begin{itemize}
\tightlist
\item
  Fanout on message write: write one message record, but create one delivery record per member.
\item
  Alternative: write one message, each member\textquotesingle s client polls for group updates using a \texttt{last\_read\_seq} cursor.
\end{itemize}

\textbf{Recommended for \(\leq\) 500 members:}

\setcodelang{}

\begin{verbatim}
Message write → Kafka
Kafka consumer → write 1 message to messages table
                → write N "unread" records to user_conversation table (one per member)
Online members → WebSocket push
Offline members → Push notification
\end{verbatim}

\textbf{Very large groups (\textgreater{} 500, e.g., broadcast channels):}

\begin{itemize}
\tightlist
\item
  Switch to fanout-on-read: message is in one place, members\textquotesingle{} clients pull new messages using \texttt{seq\_num} cursor.
\item
  Similar to news feed pull model.
\end{itemize}

\subsubsection{6.4 Message Ordering and Sequence Numbers}\label{64-message-ordering-and-sequence-numbers}

\textbf{Problem:} Two users send messages simultaneously. Who goes first?

\textbf{Solution: Sequence number per conversation (not global).}

\begin{itemize}
\tightlist
\item
  Each conversation has a monotonic \texttt{seq\_num} counter stored in Redis or a coordination service (like Zookeeper).
\item
  When Message Service processes a message, it atomically increments \texttt{seq\_num} for that \texttt{conversation\_id} and assigns it to the message.
\item
  Messages displayed in \texttt{seq\_num} order on all clients.
\end{itemize}

\textbf{Clock skew:} Client timestamps cannot be trusted. Always use server-assigned \texttt{seq\_num} for ordering. Client timestamp only shown as display time.

\setcodelang{}

\begin{verbatim}
INCR "seq:{conversation_id}"  → returns 100034 (atomic in Redis)
Store message with seq=100034
\end{verbatim}

\subsubsection{6.5 Online Presence}\label{65-online-presence}

\textbf{Heartbeat protocol:}

\begin{itemize}
\tightlist
\item
  Client sends WebSocket heartbeat every 10 seconds.
\item
  Chat Server updates Redis: \texttt{SET\ "presence:\{user\_id\}"\ 1\ EX\ 30}.
\item
  Key expires in 30s if no heartbeat \(\rightarrow\) user considered offline.
\end{itemize}

\textbf{Displaying presence:}

\begin{itemize}
\tightlist
\item
  When user A opens a chat with user B: fetch \texttt{GET\ "presence:\{user\_id\_B\}"}.
\item
  "Last seen": stored as \texttt{last\_seen:\{user\_id\}} on heartbeat stop (WebSocket disconnect or timeout).
\end{itemize}

\textbf{Scale challenge:}

\begin{itemize}
\tightlist
\item
  500M DAU \(\times\) heartbeat every 10s = 50M presence writes/sec to Redis.
\item
  \textbf{Optimization:} Only send heartbeat if the user is actively viewing a conversation. Idle app = slower heartbeat (30s). Closed app = no heartbeat.
\item
  Use presence fanout sparingly: only push presence updates to contacts who have a conversation open with that user.
\end{itemize}

\subsubsection{6.6 Read Receipts}\label{66-read-receipts}

\textbf{States:} \texttt{sent} \(\rightarrow\) \texttt{delivered} \(\rightarrow\) \texttt{read}.

\setcodelang{}

\begin{verbatim}
sent:       Message reached the chat server (ack from server to sender).
delivered:  Message delivered to recipient's device (ack from recipient device).
read:       Recipient opened the conversation (client sends "read" ack).
\end{verbatim}

\textbf{Implementation:}

\begin{itemize}
\tightlist
\item
  Each message has a \texttt{status} field in Cassandra.
\item
  Recipient sends \texttt{\{"type":"ack","msg\_id":"...","ack\_type":"read"\}} over WebSocket.
\item
  Chat Server \(\rightarrow\) Message Service \(\rightarrow\) update Cassandra status \(\rightarrow\) forward read receipt to sender\textquotesingle s Chat Server \(\rightarrow\) sender gets green double-tick.
\end{itemize}

\textbf{Group read receipts:}

\begin{itemize}
\tightlist
\item
  Store \texttt{read\_by:\ Map\textless{}user\_id,\ timestamp\textgreater{}} per message in Cassandra.
\item
  "Delivered to all" = all members\textquotesingle{} devices have acknowledged \texttt{delivered}.
\item
  Show "Read by X, Y, Z" as clients send individual read acks.
\end{itemize}

\subsection{7. Bottlenecks \& Scaling}\label{7-bottlenecks--scaling-3}

\textbf{150 million concurrent WebSocket connections:}

\begin{itemize}
\tightlist
\item
  Each Chat Server handles \textasciitilde50K--100K connections (Go or Node.js event loop).
\item
  Need \textasciitilde1,500--3,000 Chat Servers.
\item
  Chat Servers are stateful but share no business logic --- scale horizontally, no coupling.
\end{itemize}

\textbf{Kafka throughput (230K messages/sec):}

\begin{itemize}
\tightlist
\item
  Partition by \texttt{conversation\_id} (ensures ordering per conversation).
\item
  230K msg/sec \(\times\) 1 KB = 230 MB/s --- Kafka handles this easily with dozens of brokers.
\item
  Message Service consumer group: scale consumers to match Kafka partition count.
\end{itemize}

\textbf{Cassandra message storage (3 TB/day):}

\begin{itemize}
\tightlist
\item
  Partition by \texttt{conversation\_id} ensures messages in same conversation are co-located.
\item
  Hot conversation: one partition gets hammered \(\rightarrow\) add \texttt{bucket} (time shard) to partition key.
\item
  3 TB/day \(\times\) 365 = \textasciitilde1.1 PB/year. Cassandra compresses \textasciitilde3:1 \(\rightarrow\) \textasciitilde370 TB/year of disk.
\item
  TTL old messages: archive messages \textgreater{} 1 year to cold storage (S3 Glacier), keep last 1 year hot.
\end{itemize}

\textbf{Media delivery:}

\begin{itemize}
\tightlist
\item
  Don\textquotesingle t send media over WebSocket --- too large.
\item
  Client uploads directly to S3 (pre-signed URL from Media Service).
\item
  Sends media URL in the chat message. Recipients download from CDN-fronted S3.
\end{itemize}

\textbf{Push notification scale:}

\begin{itemize}
\tightlist
\item
  500M users. At any time, 70\% offline.
\item
  Surge: celebrity group sends a message \(\rightarrow\) 1M push notifications simultaneously.
\item
  Use async workers + APNs/FCM batch APIs. Rate limit per APNs connection (use multiple connections).
\end{itemize}

\subsection{8. Trade-offs Summary}\label{8-trade-offs-summary-3}

\begin{longtable}[]{@{}
  >{\raggedright\arraybackslash}p{(\columnwidth - 6\tabcolsep) * \real{0.1382}}
  >{\raggedright\arraybackslash}p{(\columnwidth - 6\tabcolsep) * \real{0.2910}}
  >{\raggedright\arraybackslash}p{(\columnwidth - 6\tabcolsep) * \real{0.2110}}
  >{\raggedright\arraybackslash}p{(\columnwidth - 6\tabcolsep) * \real{0.3598}}@{}}
\toprule\noalign{}
\rowcolor{acc!16}
\begin{minipage}[b]{\linewidth}\raggedright
Decision
\end{minipage} & \begin{minipage}[b]{\linewidth}\raggedright
Choice Made
\end{minipage} & \begin{minipage}[b]{\linewidth}\raggedright
Alternative
\end{minipage} & \begin{minipage}[b]{\linewidth}\raggedright
Reason
\end{minipage} \\
\midrule\noalign{}
\endhead
\bottomrule\noalign{}
\endlastfoot
Real-time transport & \textbf{WebSocket} & Long polling / SSE & Full-duplex; lowest latency; widely supported \\
\rowcolor{zebra}
Message storage & \textbf{Cassandra (conv-partitioned)} & MySQL (sharded) & Write throughput; time-series access pattern; no joins needed \\
Connection registry & \textbf{Redis (TTL-based)} & Consistent hash routing & Fast lookup; auto-evicts stale connections on TTL expiry \\
\rowcolor{zebra}
Message ordering & \textbf{Server-assigned seq\_num (Redis INCR)} & Vector clocks & Simple; sufficient for single-conversation ordering \\
Group delivery & \textbf{Fanout to members (\textless{} 500) / pull (\textgreater{} 500)} & Always pull & Balances latency vs write amplification by group size \\
\rowcolor{zebra}
Presence heartbeat & \textbf{10s client heartbeat, 30s TTL} & Server-push disconnect events & Handles silent disconnects (network failure, app killed) \\
Read receipts & \textbf{Per-message status in Cassandra} & Separate receipt table & Avoids join; co-located with message data \\
\rowcolor{zebra}
Media & \textbf{S3 + CDN (URL in message)} & In-message blob & Media over WebSocket is impractical at scale \\
\end{longtable}

\emph{End of Part 1 --- Foundational System Design Problems.}

\clearpage
\section{Visual Reference}
\noindent A set of figures distilling the key quantitative intuitions of this topic.\par\vspace{4pt}
\visualfigure{../figures/sdp-1/01-token-bucket.pdf}{A bucket refills at a steady rate up to a capacity. Each request spends a token; an empty bucket throttles. This smooths traffic while allowing short bursts.}
\end{document}
